\section{Universal constructions}\label{sec:universal_constructions}

\paragraph{Representable functors}

\begin{definition}\label{def:functor_representation}
  A \term{representation} of the functor \( F: \cat{C} \to \cat{Set} \) is a pair \( (A, \rho) \), where \( A \) is an object in \( \cat{C} \) and \( \rho: \hom(\anon*, A) \to F \) is a \hyperref[def:natural_isomorphism]{natural isomorphism}.

  We call \( F \) \term{representable} if a representation exists.
\end{definition}

\begin{proposition}\label{thm:def:functor_representation}
  \hyperref[def:functor_representation]{Functor representations} have the following basic properties:
  \begin{thmenum}
    \thmitem{thm:def:functor_representation/uniqueness} All representing objects for the same functor are isomorphic.
  \end{thmenum}
\end{proposition}
\begin{proof}
  \SubProofOf{thm:def:functor_representation/uniqueness} Suppose we have two natural isomorphisms
  \begin{align*}
    &\rho: \hom(\anon*, A) \to F, \\
    &\sigma: \hom(\anon*, B) \to F.
  \end{align*}

  Then, by \cref{thm:def:morphism_invertibility/composition}, the transformation \( \sigma \bincirc \rho^{-1}: \hom(\anon*, A) \to \hom(\anon*, B) \) is also an isomorphism.

  \Cref{thm:isomorphisms_via_yoneda_embedding} implies that \( A \) and \( B \) are isomorphic.
\end{proof}
